\section{Discussion}\label{sec:disc}

\subsection{Relation to earlier work}\label{ss:disc-prior}
Hedden, Herald and Kirk verify their Conjecture~6.5 (Conjecture~\ref{conj:hhk65}) for two-bridge knots \cite[\S7]{HHK2} and for a
list of torus knots \cite[\S11]{HHK2}; the torus knots $T(3,n)$ are treated in \cite{WuebbenV}.  For
a decomposition of their own, Fukumoto, Kirk and Pinz\'on-Caicedo describe the unperturbed traceless
character variety of the pretzel knots and count nine generators for $P(-2,3,7)$ \cite[\S5]{FKP}; they do not perturb or compute a Floer complex.  Theorems~\ref{thm:hhk}
and~\ref{thm:twisted} supply the perturbed curves, their Floer complexes and the gradings.  For
$q\ge9$ the differential in Theorem~\ref{thm:hhk} is nonzero, so the agreement with $\Inat$ is not a
count of generators.  The twists act on the pillowcase by the shear of \cite[(28)--(30)]{HHK2}.

For Smith's decomposition of $K_q$ as a sum of rational tangles, Smith proves that for $q=5$ the pairing with
zero bounding cochains has rank $5$ rather than $\dim\Inat=7$~\cite[proof of Thm.~5.9]{Smith}, so that under the
conjecture of Cazassus, Herald, Kirk and Kotelskiy a nonzero bounding cochain is required
there~\cite[Thm.~5.9]{Smith}. The decomposition used here needs none. Maurer--Cartan elements for Smith's decomposition at
$q=5,7$, and their relation to the Bar-Natan invariant of Kotelskiy, Watson and Zibrowius~\cite{KWZ} of the tangle formed
by the twist regions with $3$ and $q$ half-twists, are the subject of~\cite{WuebbenIIb}.

\subsection{Open problems}\label{ss:disc-open}
\emph{Chain level.} The pillowcase complex of Theorem~\ref{thm:hhk} is an S-complex over $\F_2$ whose only nonzero
component is the differential of its irreducible part (Remark~\ref{rem:scomplexP}). Whether it is S-chain homotopy
equivalent to $\tC(K_q;\F_2)$ is open. For the torus knots $T(3,n)$ the corresponding statement
holds~\cite[Thm.~1.3]{WuebbenV}, and there the pillowcase differential is a component of type $\delta_1$; for $K_q$
with $q\ge9$ it lies in the irreducible part instead.

\emph{Other families.}  Theorem~\ref{thm:twisted} covers $T(3,n;2,m)$ for $n\equiv1,2\pmod6$.  For
$n\equiv4,5\pmod6$ and $n\ge10$ the curve of \cite{WuebbenV} has a central circle, and the shear
argument of Section~\ref{sec:sheared} needs that circle to be a straight line, which we prove only for
$n=5$; the signature argument also fails for $n\equiv5\pmod6$.
